% =============================================================
% STYLE GUIDE FOR THE MSc IN PHYSICS (FUNDAMENTAL PHYSICS)
% Compile it from inside the style_guide/ folder:
%pdflatex style_guide -> bibtex style_guide -> pdflatex style_guide -> pdflatex style_guide
% It reuses thesis.cls, thesis.bib, aa_url.bst and images/ from
% the template folder (one level up).
% =============================================================

% Look for the class, style and input files in the parent folder too
\makeatletter
\def\input@path{{../}}
\makeatother

\documentclass[a4paper,10pt,openany]{thesis}
\usepackage{silence}
\WarningsOff[latex/command/changed]

\graphicspath{{../}}               % figures are taken from ../images/
\usepackage{epstopdf}
\epstopdfsetup{outdir=./}          % write converted EPS files here, not in ../
\usepackage{enumitem}
\setlist{nosep,leftmargin=*}
\usepackage{capt-of}

% Compact layout for this guide only (the thesis itself uses 1.5 spacing)
\renewcommand{\baselinestretch}{1.0}
\setlength{\textheight}{26cm}
\setlength{\topmargin}{-2cm}
\setlength{\parskip}{3pt}
\setlength{\parindent}{0pt}
\renewcommand{\thesection}{\arabic{section}}
\renewcommand{\thefigure}{\arabic{figure}}
\renewcommand{\thetable}{\arabic{table}}
\renewcommand{\theequation}{\arabic{equation}}
\pagestyle{plain}
\makeatletter
\def\verbatim@font{\small\ttfamily}  % smaller font for code examples
\makeatother

% Yachay Tech logo in the upper right corner of every page
\AddToHook{shipout/background}{%
    \put(\paperwidth-4.2cm,-1.6cm){\includegraphics[width=3cm]{images/logos/yachay_logo.pdf}}}

\begin{document}


\begin{center}
    {\Large\bfseries MSc in Physics, specialisation in Fundamental Physics\\
    Yachay Tech University\\[5pt]
    Thesis style guide:}\\[2pt]
    
\end{center}

This guide summarises the style and shows how to use the \LaTeX{} template. If something here differs from the official documents (\href{https://yachaytecheduec-my.sharepoint.com/:w:/r/personal/posgrados_yachaytech_edu_ec/Documents/4.5.2\%20POSGRADOS/FORMATOS/Estudiantes/Formatos\%20para\%20titulaci\%C3\%B3n/Anexo\%203.\%20Formato\%20de\%20presentaci\%C3\%B3n.docx?d=w914bf4407dba4fd4ab353bf88fc8fdef&csf=1&web=1&e=EQSmdg}{Anexo 3. Formato de presentación} and \href{https://www.yachaytech.edu.ec/wp-content/uploads/2021/11/RCIFI-S0-011-No.-0320-2021-Reglamento-integral-de-postgrados.pdf}{Reglamento Integral de Postgrados}), the official documents apply.

\section{Language}
\begin{itemize}
     \item Write the thesis in English. Your oral defence will also be in English. Pick British or American spelling and be consistent throughout the thesis.
    \item The cover page (carátula) is in Spanish, but the thesis title on it is in English, after the word ``Título:''. The Autoría and Autorización de publicación pages are in Spanish. All of these are generated by the template; you only fill in your details in \texttt{main.tex}.
    \item Include a \textit{Resumen} (Spanish) and an \textit{Abstract} (English) of about 200-250 words each (not longer, please), each one on its own page and followed by 3 to 6 keywords.
\end{itemize}

\section{Format and order}
\begin{itemize}
    \item The template already sets the format: A4 paper, Times font, 10\,pt, 1.5 line spacing. Do not change it, and do not add page headers.
    \item Pages are numbered from the table of contents onwards. The cover, Autoría, Autorización de publicación and initial sections have no page number.
    \item Order: cover page; Autoría and Autorización de publicación (electronically signed); Dedicatoria and Agradecimiento (optional, one page each); Resumen; Abstract; table of contents; chapters; references; appendices.
\end{itemize}

\section{Length and content}
\begin{itemize}
    \item Aim for 50 to 70 pages in total, and no more than 100 pages including everything (figures and appendices). The main text (chapters only) is usually 40 to 60 pages, or about 12\,000 to 18\,000 words.
    \item Quality matters more than length: 40 clear pages with your own results are better than a 100 pages that repeat textbook material. A good target is a thesis that could become a journal article.
    \item The length depends on the topic: theoretical work with long derivations, numerical work with simulations, and experimental work with set-ups and error analysis all need different spaces. Discuss with your supervisors the content distribution.
    \item The thesis must show independent research, a critical review of the literature, and original results, with the methods, results and conclusions clearly explained.
    \item If parts of your thesis are published, submitted or on arXiv, say so: uncomment the \textit{Disclaimer} section in \texttt{main.tex} and list the papers with their links (DOI or arXiv).
    \item The chapter titles in the template are only a suggestion; you can rename, merge or split chapters. A typical split is: Introduction (5--8 pages), Theoretical background (8--12), Methodology (8--12), Results and discussion (16--24), Conclusions (3--4).
    \item The similarity index from the anti-plagiarism software (Yachay Tech uses Compilatio: \href{https://www.compilatio.net}{www.compilatio.net}) must not exceed 10\% (\href{https://www.yachaytech.edu.ec/wp-content/uploads/2021/11/RCIFI-S0-011-No.-0320-2021-Reglamento-integral-de-postgrados.pdf}{Reglamento Integral de Postgrados}, Art.~78). Please avoid pasting AI-generated text; Compilatio will flag the LLM-generated text.
\end{itemize}

%    \section{Binding and CD}
%    Black hard covers with gold letters, A4. On the spine, in gold: the graduation year at the bottom, then your name and the title, and \textbf{ECFN} with the School logo at the top. Add a CD in an envelope glued inside the back cover, with the thesis as a PDF named \textit{TRABAJO DE TITULACIÓN -- APELLIDO Y NOMBRE DEL AUTOR}. Deliver both to the Library. See Anexo 4 for the layouts.

\section{Using the template}

\textbf{Filling out the template.} Edit only these files:
\begin{itemize}
    \item \texttt{main.tex}: the main file, the one you compile. Fill in your details for the cover (title in English after ``Título:'', name, ID number, supervisor, month in Spanish and year), and write the Dedicatoria, Agradecimiento, Resumen and Abstract. The citation style is APA-like by default; to change it (e.g.\ to the astrophysics style), follow the comment at the end of \texttt{main.tex}.
    \item \texttt{content/}: one file per chapter (\texttt{01\_introduction.tex}, \ldots) and one per appendix. Write your thesis here. To rename, add or remove a chapter, edit the file and its \verb|\input| line in \texttt{main.tex}.
    \item \texttt{images/plots/}: your figures, preferably as PDF or EPS.
    \item \texttt{thesis.bib}: your references in BibTeX format (journals, arXiv, NASA ADS and INSPIRE-HEP export it). Some styles change titles to lower case; put words that must keep their capitals in braces, e.g.\ \verb|{Klein-Gordon}|.
    \item Do not edit \texttt{thesis.cls} (format and Spanish pages) or \texttt{aa\_url.bst} (astrophysics citation style).
\end{itemize}
Compile \texttt{main.tex} with pdflatex, bibtex, pdflatex, pdflatex (Overleaf and \texttt{latexmk} do this automatically).

\textbf{Example.} The text below uses a figure, a table and three citations:

\begin{quote}
The function $\sin x$ is shown in Figure~\ref{fig:sine}, and some of its values are listed in Table~\ref{tab:example}. Examples of articles in different areas are on quantum scattering \citep{PhysRevA.71.052101}, surfaces \citep{PhysRevB.70.125402} and social dynamics \citep{PhysRevE.110.024311}. \citet{PhysRevA.71.052101} studied a Klein-Gordon particle.
\end{quote}
\vspace{-4pt}
\begin{verbatim}
The function $\sin x$ is shown in Figure~\ref{fig:sine}, and some
of its values are listed in Table~\ref{tab:example}. ... on quantum
scattering \citep{PhysRevA.71.052101}, ... \citet{PhysRevA.71.052101}
\end{verbatim}
Use \verb|\citep{key}| for citations in parentheses and \verb|\citet{key}| when the authors are part of the sentence.

\textbf{Figures.} Save them in \texttt{images/plots/}. For the final version, use vector plots (PDF or EPS) whenever possible; for images that cannot be vector (photos, maps, snapshots of simulations), use a high resolution of at least 300 dpi (PNG or JPG). Refer to each one in the text (Figure~\ref{fig:sine}) and give it a caption below. With \texttt{overpic} you can write labels on top of an image (\verb|\put(x,y){...}|, with $x$ and $y$ in percent of the figure width), so you can change them later without redoing the plot.

\begin{minipage}[t]{0.55\textwidth}
\vspace{0pt}
\begin{verbatim}
\begin{figure}[H]
  \centering
  \begin{overpic}[width=0.5\textwidth]
      {images/plots/sine.eps}
    \put(55,-3){$x$ [rad]}
    \put(36,50){$f(x)=\sin x$}
  \end{overpic}
  \caption{The function $\sin x$.}
  \label{fig:sine}
\end{figure}
\end{verbatim}
\end{minipage}%
\begin{minipage}[t]{0.45\textwidth}
\vspace{0pt}
\centering
    \begin{overpic}[width=0.95\textwidth]{images/plots/sine.eps}
        \put(52,-4){\small $x$ [rad]}
        \put(36,49){\small $f(x)=\sin x$}
    \end{overpic}
    \vspace{2pt}
    \captionof{figure}{The function $\sin x$.}
    \label{fig:sine}
\end{minipage}

\medskip
\textbf{Tables.} Put the caption above the table, use horizontal lines only, and give units in the headers (Table~\ref{tab:example}).

\begin{minipage}[t]{0.55\textwidth}
\vspace{0pt}
\begin{verbatim}
\begin{table}[H]
  \centering
  \caption{Values of $\sin x$.}
  \label{tab:example}
  \begin{tabular}{cc}
    \hline $x$ [rad] & $\sin x$ \\ \hline
    $0$ & $0$ \\ $\pi/2$ & $1$ \\ \hline
  \end{tabular}
\end{table}
\end{verbatim}
\end{minipage}%
\begin{minipage}[t]{0.45\textwidth}
\vspace{0pt}
\centering
\captionof{table}{Values of $\sin x$.}
\label{tab:example}
\begin{tabular}{cc}
\hline
$x$ [rad] & $\sin x$ \\
\hline
$0$     & $0$ \\
$\pi/2$ & $1$ \\
\hline
\end{tabular}
\end{minipage}

\medskip
\textbf{Equations.} Number the equations you refer to, give them a label, and refer to them with \verb|\eqref|, e.g.\ the 1D Schrödinger equation, Eq.~\eqref{eq:schrodinger}. End equations with a comma or full stop when the sentence needs it.

\begin{minipage}[t]{0.55\textwidth}
\vspace{0pt}
\begin{verbatim}
\begin{equation}
  i\hbar\frac{\partial\psi}{\partial t}
    = -\frac{\hbar^2}{2m}
       \frac{\partial^2\psi}{\partial x^2}
     + V(x)\,\psi ,
  \label{eq:schrodinger}
\end{equation}
\end{verbatim}
\end{minipage}%
\begin{minipage}[t]{0.45\textwidth}
\vspace{-4pt}
\begin{equation}
  i\hbar\frac{\partial\psi}{\partial t} = -\frac{\hbar^2}{2m}\frac{\partial^2\psi}{\partial x^2} + V(x)\,\psi ,
  \label{eq:schrodinger}
\end{equation}
\end{minipage}

% References printed as a section (not a new page)
\begingroup
\makeatletter
\renewcommand{\chapter}{\@ifstar{\section*}{\section}}
\makeatother
\footnotesize
\bibliography{../thesis}
\bibliographystyle{apalike}
\endgroup

\end{document}
